% =================================================
% Решения домашней работы. Урок 1
% =================================================

\clearpage
\homeworktitle{Решения и комментарии}{Версия для учителя}

\begin{teacheranswer}[Задание 1]
  \begin{itemize}
    \item \(36\): две цифры --- \(3,6\);
    \item \(507\): три цифры --- \(5,0,7\);
    \item \(70\,090\): пять цифр --- \(7,0,0,9,0\).
  \end{itemize}
\end{teacheranswer}

\begin{teacheranswer}[Задание 2]
  \begin{enumerate}
    \item цифра сотен числа \(48\,372\) --- \(3\);
    \item цифра десятков тысяч числа \(605\,814\) --- \(0\);
    \item цифра десятков числа \(90\,407\) --- \(0\).
  \end{enumerate}
\end{teacheranswer}

\begin{criterionbox}
  Нулевые ответы во втором и третьем пунктах специально проверяют, считает
  ли ученик ноль полноценной цифрой позиционной записи.
\end{criterionbox}

\begin{teacheranswer}[Задание 3]
  \begin{enumerate}
    \item \(7\,000\);
    \item \(3\,000\);
    \item \(90\).
  \end{enumerate}
\end{teacheranswer}

\begin{teacheranswer}[Задание 4]
  \begin{align*}
    64\,205&=60\,000+4\,000+200+5,\\
    8\,090&=8\,000+90,\\
    300\,407&=300\,000+400+7.
  \end{align*}
\end{teacheranswer}

\begin{teacheranswer}[Задание 5]
  \begin{enumerate}
    \item \(56\,302\);
    \item \(90\,740\);
    \item \(425\,009\).
  \end{enumerate}
\end{teacheranswer}

\begin{teacheranswer}[Задание 6]
  Цифра \(3\) в числе \(71\,304\) стоит в разряде сотен, а не тысяч.
  Поэтому:
  \[
    71\,304=70\,000+1\,000+300+4.
  \]
\end{teacheranswer}

\begin{teacheranswer}[Задание 7*]
  \[
    6\cdot10\,000+34\cdot100+8
    =60\,000+3\,400+8
    =\boxed{63\,408}.
  \]
\end{teacheranswer}

\begin{resultbox}[Краткие ответы]
  \begin{enumerate}
    \item \(2;\ 3;\ 5\) цифр.
    \item \(3;\ 0;\ 0\).
    \item \(7\,000;\ 3\,000;\ 90\).
    \item \(60\,000+4\,000+200+5;\ 8\,000+90;\ 300\,000+400+7\).
    \item \(56\,302;\ 90\,740;\ 425\,009\).
    \item \(70\,000+1\,000+300+4\).
    \item \(63\,408\).
  \end{enumerate}
\end{resultbox}
